\documentclass[11pt,a4paper]{article}
\usepackage[T1]{fontenc}
\usepackage[utf8]{inputenc}
\usepackage{lmodern}
\usepackage{amsmath,amssymb,amsthm}
\usepackage{booktabs}
\usepackage{graphicx}
\usepackage[margin=2.5cm]{geometry}
\usepackage{hyperref}
\hypersetup{colorlinks=true,linkcolor=blue,citecolor=blue,urlcolor=blue}
\graphicspath{{../figures/}}

\newtheorem{proposition}{Proposition}
\newtheorem{lemma}{Lemma}
\newtheorem{assumption}{Assumption}
\theoremstyle{remark}
\newtheorem{remark}{Remark}

\newcommand{\wrap}{\operatorname{wrap}}
\newcommand{\med}{\operatorname{med}}
\newcommand{\E}{\mathbb{E}}
\newcommand{\Var}{\operatorname{Var}}
\newcommand{\Cov}{\operatorname{Cov}}
\newcommand{\sgn}{\operatorname{sgn}}
\newcommand{\dg}{^{\circ}}
\newcommand{\Wt}{W_{\mathrm{t}}}
\newcommand{\Nt}{N_{\mathrm{t}}}
\newcommand{\Wr}{W_{\mathrm{r}}}
\newcommand{\Nr}{N_{\mathrm{r}}}

\title{Static yaw misalignment from high-resolution SCADA:\\
       derivation of the estimators C1--C5, the consensus residual\\
       and the submission model}
\author{Nadara / WeDoWind challenge, participant 33}
\date{September 2026}

\begin{document}
\maketitle

\begin{abstract}
This note derives, from first principles, every quantity used in the pipeline
\texttt{src/yaw/}: the sensor model that links the SCADA direction channels to the
physical yaw error (\S\ref{sec:sensor}), the aerodynamic power--yaw law
(\S\ref{sec:aero}), the four power-based estimators of the static misalignment
(C1--C4, \S\ref{sec:c1}--\S\ref{sec:c4}), the vane-gain idea C5 and why the data
encoding forbids it (\S\ref{sec:c5}), the neighbour-consensus heading residual and
its change-point segmentation (\S\ref{sec:consensus}), the submission model that
combines them (\S\ref{sec:submission}), and the scoring rule against which all of
this is judged (\S\ref{sec:scoring}). Each section states the assumptions
explicitly, derives the estimator, gives its sampling variance, and identifies the
systematic error that dominates it in practice. The measured outcomes
(\texttt{docs/03\_estimators.md}, \texttt{docs/04\_submission.md}) are quoted where
they decide between competing models.
\end{abstract}

\tableofcontents

%======================================================================
\section{Notation and the sensor model}\label{sec:sensor}

\subsection{Angles}
All directions are compass bearings in degrees, clockwise from north, and every
difference of two bearings is wrapped to $(-180\dg,180\dg]$,
\begin{equation}
  \wrap(x) = \big((x+180) \bmod 360\big) - 180 .
\end{equation}
For a set of bearings $x_1,\dots,x_n$ the \emph{circular mean} is
\begin{equation}
  \bar x = \operatorname{atan2}\Big(\tfrac1n\sum_i \sin x_i,\ \tfrac1n\sum_i \cos x_i\Big),
  \qquad
  \bar R = \Big|\tfrac1n\sum_i e^{\mathrm{i}x_i}\Big| \in [0,1],
  \label{eq:circmean}
\end{equation}
where the resultant length $\bar R$ measures concentration ($\bar R=1$: all
equal). The \emph{circular median} used throughout the consensus code is the
ordinary median of the wrapped deviations from the circular mean,
$\med_\circ(x) = \bar x + \med_i \wrap(x_i - \bar x)$; it inherits the
robustness of the median while being invariant to where the $\pm180\dg$ cut sits.

\subsection{True and reported quantities}
Let, at a given instant,
\begin{align*}
  \Wt &\ \text{true (free-stream) wind direction at the rotor,}\\
  \Nt &\ \text{true nacelle heading (rotor axis),}\\
  \gamma &= \wrap(\Wt - \Nt) \quad \text{the instantaneous yaw error,}
\end{align*}
with the sign convention that $\gamma>0$ when the wind comes from clockwise of
the rotor axis. The SCADA stream does not report either true quantity. It reports
\begin{align}
  \Nr &= \Nt + \delta_n, \label{eq:Nr}\\
  v_m &= b\,\gamma + \delta_v, \label{eq:vm}\\
  \Wr &= \Nr + v_m, \label{eq:Wr}
\end{align}
where $\delta_n$ is the nacelle encoder's mounting offset (the encoder has no
north reference), $\delta_v$ is the wind-vane's mounting offset and $b\le1$ is the
vane's gain. Equation~\eqref{eq:Wr} is not an assumption but a fact of this data
set: the field \texttt{WindDir} is computed by the controller as heading plus
vane, so that \texttt{WindDir}$-$\texttt{NacDir} returns the vane reading exactly
(verified in \texttt{docs/01}). The gain $b<1$ models the well-documented
damping of the vane signal behind a rotating rotor: the wake swirl and the
nacelle body deflect the flow so that the vane sees a fraction of the true inflow
angle.\footnote{See the literature on nacelle transfer functions, e.g.\ the
2025 WES preprint cited in \texttt{PLAN.md}. Everything below is written for
general $b$; the implementation uses $b=1$ wherever no independent estimate
exists.}

\subsection{The controller and the definition of the label}
The yaw controller integrates $v_m$ over minutes and rotates the nacelle whenever
the filtered reading leaves a deadband around its set-point $s$. In steady state
the controller therefore enforces $\E[v_m]=s$, so from~\eqref{eq:vm}
\begin{equation}
  \theta \;\equiv\; \E[\gamma] \;=\; \frac{s-\delta_v}{b}.
  \label{eq:theta_def}
\end{equation}
$\theta$ is the \emph{static yaw misalignment}: the persistent yaw error the
controller cannot see because it is inside its own measurement. It is what the
WindFit reference system measures with an independent sensor on the nacelle,
and it is the label of the challenge.\footnote{WindFit fits the power-optimal
heading over several days; its daily series is therefore smoothed and drifts for
a week before each step, which is why days near a step are excluded from
scoring (\S\ref{sec:scoring}).} The label is negative on all three training
turbines, and the empirical sign check of \S\ref{sec:signcheck} fixes that the
label sign coincides with $\gamma$ as defined here.

Two consequences shape everything that follows.
\begin{enumerate}
\item \textbf{The vane cannot see $\theta$.} Averaged over any window the vane
      reads $s$, whatever $\theta$ is. Measured: the daily median of $v_m$ is
      $\approx-3\dg$ on every PPP turbine irrespective of the label
      (correlation with the label $-0.33$). Information about $\theta$ must come
      from the \emph{response} of the turbine to $\gamma$ (power, rotor speed)
      or from an external direction reference (neighbours).
\item \textbf{$\theta$ changes in three ways}, and they are not equivalent for the
      sensors: a change of the set-point $s$ (controller parameter), a change of
      $\delta_v$ (vane re-mounted or re-calibrated), or a change of $b$
      (blade/rotor state). All three move $\theta$ by~\eqref{eq:theta_def}, but
      only the first two move the neighbour residual of \S\ref{sec:consensus},
      and an encoder re-reference ($\delta_n$) moves the residual without moving
      $\theta$ at all.
\end{enumerate}

%======================================================================
\section{Aerodynamic model of power under yaw}\label{sec:aero}

\subsection{The $\cos^k$ law}
Actuator-disc theory for a rotor yawed by $\gamma$ replaces the free-stream speed
$U$ by its component along the rotor axis, $U\cos\gamma$, in the mass flow and in
the energy extracted. With $P=\tfrac12\rho A C_p U^3$ for the aligned rotor this
gives $P(\gamma) = P(0)\cos^3\gamma$. Measurements on real rotors give a smaller
exponent because the swept disc also captures part of the tangential component
and because the tip-speed ratio, and hence $C_p$, shifts with $\gamma$ under
torque control. The standard empirical form is
\begin{equation}
  P(U,\gamma) = \tfrac12\rho A\, C_p(\lambda,\beta)\, U^3 \cos^k\gamma ,
  \qquad k\in[1.4,\,3],
  \label{eq:cosk}
\end{equation}
which the OpenOA static-yaw method also assumes. The implementation fixes $k=2$
after finding that a free $k$ wanders between $0.2$ and $13$ from month to month
on this data (\S\ref{sec:c1_bias}); with $k$ fixed the only unknown of the
power--yaw curve is its location.

\subsection{Small-angle behaviour and why the deadband matters}
Expanding \eqref{eq:cosk} around $\gamma=0$ with $\gamma$ in radians,
\begin{equation}
  \cos^k\gamma = 1 - \frac{k}{2}\gamma^2 + O(\gamma^4),
  \qquad
  \ln\cos^k\gamma = -\frac{k}{2}\gamma^2 - \frac{k}{12}\gamma^4 + O(\gamma^6).
  \label{eq:cosk_small}
\end{equation}
The controller keeps $v_m$ inside roughly $\pm8\dg$ (measured occupancy in
\texttt{figures/03\_estimators.png}). Over that range, with $k=2$, the power
varies by $\tfrac{k}{2}(8\pi/180)^2 \approx 2\,\%$. Every estimator that locates
the peak of \eqref{eq:cosk} from operational data therefore has to resolve a
curvature whose full effect is two per cent of power, against a background of
sector-, turbulence- and season-dependent power variations of the same order.
This is the root cause of the failures documented for C1 and C2.

\subsection{Nacelle anemometer under yaw}
The nacelle anemometer sits behind the rotor and reads
\begin{equation}
  U_n = U\, h(\gamma,\lambda), \qquad h(0,\lambda)=h_0(\lambda)<1,\quad
  \partial_\gamma h \ne 0 ,
  \label{eq:anemo}
\end{equation}
i.e.\ the rotor-induced deficit itself depends on the yaw angle. Binning by
$U_n$ (C1) is therefore \emph{not} binning by $U$: at fixed $U_n$ a yawed rotor
has a different free-stream speed than an aligned one, and the apparent
power--yaw curve is $\cos^k\gamma / h^3(\gamma,\lambda)$ rather than
$\cos^k\gamma$. Since $h$ is symmetric in $\gamma$ to leading order this does not
move the peak by itself, but it flattens or sharpens the curve, i.e.\ it is
absorbed into an effective $k$ --- another reason to fix $k$.

%======================================================================
\section{C1: power-versus-vane peak with a nacelle wind reference}\label{sec:c1}

\subsection{Estimator}
Take 1-minute means in partial load, on healthy days, with the wind outside the
wake sectors of the layout (\S\ref{sec:wake}), and with a within-minute vane
standard deviation below $6\dg$ so that the minute has a well-defined yaw angle.
Let $\bar v$ be the minute's mean vane and $\bar P$ its mean power. The
free-stream speed is unknown, so it is removed by \emph{normalisation within
wind-speed bins}: with $B(u)$ the $0.5\,m/s$ bin of the nacelle wind speed,
\begin{equation}
  \tilde P = \frac{\bar P}{\med\{\bar P' : \bar P' \in B(u)\}} .
  \label{eq:norm}
\end{equation}
By \eqref{eq:cosk}, and because the bin median is taken over all yaw angles in
the bin (which are concentrated near the set-point $s$),
\begin{equation}
  \E[\tilde P \mid \bar v] \;\approx\; A\,\cos^k\!\big(\bar v - \theta_0\big),
  \qquad
  \theta_0 = \delta_v \ \ (\text{for } b=1),
  \label{eq:c1_model}
\end{equation}
where the peak location $\theta_0$ is the vane reading at which the rotor is
aligned, $\gamma=0 \Leftrightarrow v_m=\delta_v$. For $b\ne1$ the curve becomes
$\cos^k((\bar v-\delta_v)/b)$: the peak is still at $\delta_v$ but the curvature
is stretched by $1/b^2$, which with fixed $k$ inflates the fitted
$A$ and leaves $\theta_0$ unbiased to first order.

The data are binned by $\bar v$ in $1\dg$ steps, each bin represented by its
median $y_j$ and count $n_j$, and the two parameters $(A,\theta_0)$ are obtained
by weighted non-linear least squares,
\begin{equation}
  (\hat A,\hat\theta_0) = \arg\min_{A,\theta_0}\ \sum_j n_j\big(A\cos^k(v_j-\theta_0) - y_j\big)^2 ,
  \label{eq:c1_ls}
\end{equation}
solved with a bounded trust-region method. The bin median rather than mean is
used because the within-bin distribution of $\tilde P$ is heavy-tailed (start-up
minutes, curtailment tails that the power window does not fully remove).

\subsection{From $\theta_0$ to the misalignment}
Combining \eqref{eq:c1_model} with the controller identity
\eqref{eq:theta_def},
\begin{equation}
  \hat\theta = \frac{s - \hat\theta_0}{b} \approx \med(v_m) - \hat\theta_0 ,
  \label{eq:c1_theta}
\end{equation}
i.e.\ the misalignment is the distance between where the controller \emph{holds}
the vane and where the power \emph{peaks} on the vane axis. Equation
\eqref{eq:c1_theta} is what turns a peak location into a label estimate; the
literature often reports $-\hat\theta_0$ alone, which silently assumes $s=0$.

\subsection{Sampling variance}
Linearising the residual vector $r(p)$ of \eqref{eq:c1_ls} at the solution with
Jacobian $J$ gives the Gauss--Newton covariance
\begin{equation}
  \widehat{\Cov}(\hat p) = \hat\sigma^2 (J^\top J)^{-1},
  \qquad
  \hat\sigma^2 = \frac{\|r(\hat p)\|^2}{m - \dim p},
  \label{eq:gn_cov}
\end{equation}
with $m$ the number of bins; the reported \texttt{theta\_se} is the square root
of its $\theta_0$ diagonal element. This is a conditional standard error: it
assumes the bin medians scatter independently around a $\cos^k$ curve. The
measured quarter-to-quarter scatter of $\hat\theta_0$ within a single label state
($4$--$9\dg$) is five to ten times larger than \eqref{eq:gn_cov}, which proves
the model error dominates.

\subsection{Why the peak location is biased: the tilt lemma}\label{sec:c1_bias}
\begin{lemma}[Peak shift under a linear nuisance]\label{lem:tilt}
Let the true bin curve be $y(v)=A[1-\tfrac{k}{2}\rho^2(v-\theta_0)^2]$ with
$\rho=\pi/180$ (degrees to radians), and let an unmodelled nuisance add a
linear trend $\varepsilon\,v$ (in units of normalised power per degree). Then the
peak of the observed curve is at
\begin{equation}
  v^\ast = \theta_0 + \frac{\varepsilon}{A\,k\,\rho^2}.
\end{equation}
\end{lemma}
\begin{proof}
$\partial_v\big[y(v)+\varepsilon v\big] = -A k\rho^2 (v-\theta_0) + \varepsilon = 0$.
\end{proof}
With $k=2$, $\rho^2 = 3.05\times10^{-4}$, $A\approx1$: a tilt of one per cent of
power across the $8\dg$ deadband, $\varepsilon = 0.01/8 = 1.25\times10^{-3}$
per degree, shifts the peak by
\[
  \frac{1.25\times10^{-3}}{2\times3.05\times10^{-4}} \approx 2\dg .
\]
Any effect that correlates power with the sign of the vane deviation at the
one-per-cent level --- a wind-direction sector with slightly different shear or
turbulence, or the simple fact that the vane deviates to the right when the wind
is veering, and veering periods have their own power statistics --- moves the
fitted misalignment by degrees. The measured curves change shape between
quarters of the same label state (monotone rising, monotone falling, peaked;
\texttt{figures/03\_estimators.png}), which is exactly this lemma at work. The
same mechanism explains why the free-$k$ fit is unstable: $k$ and the tilt are
nearly collinear on a $\pm8\dg$ support.

This is the failure the challenge README warns about (``ranks the three training
turbines in the wrong order'') and the one the leading team's T0 entry
reproduces (bias $-3\dg$).

%======================================================================
\section{C2: neighbour power as the wind reference}\label{sec:c2}

C2 keeps the model \eqref{eq:c1_model}--\eqref{eq:c1_ls} but replaces the
nacelle wind-speed bin in \eqref{eq:norm} by a bin of the \emph{simultaneous
power of the nearest healthy neighbour} $P_j$ (bins of 150 power units, both
turbines in partial load and unwaked):
\begin{equation}
  \tilde P_{\mathrm{C2}} = \frac{\bar P_i}{\med\{\bar P_i' : P_j' \in B(P_j)\}} .
\end{equation}
The motivation is \eqref{eq:anemo}: the neighbour's power is a wind reference
that does not sit in the target's wake and is not a function of the target's
yaw angle. Under the assumptions
\begin{assumption}
(i) the neighbour's free-stream speed is a monotone function of the target's,
$U_j = g(U_i)$ with $g$ stable over the window; (ii) the neighbour's own yaw
error is stationary over the window,
\end{assumption}
\noindent the ratio has the same peak at $\theta_0=\delta_v$ and the derivation
of C1 goes through unchanged.

What C2 removes and what it adds is decided by Lemma~\ref{lem:tilt}. It removes
the yaw-dependence of the anemometer, which is symmetric in $\gamma$ and does
not shift the peak in any case; it adds the neighbour's sector dependence: the
map $g$ depends on the wind direction (terrain, partial wakes not in the
layout mask), and wind direction correlates with the vane deviation. Empirically
the C2 standard errors are $3$--$6\dg$ and its levels ($-5$, $-18$, $-18\dg$ on
the three training turbines) are further from the labels than C1. The neighbour
reference is the right idea for the \emph{speed} bias but the wrong tool for the
\emph{tilt} bias.

%======================================================================
\section{C3: the rotor-speed method}\label{sec:c3}

C3 was planned (Castellani et al.) and not implemented; the derivation shows why
it cannot add information on this data. In the partial-load region the
generator torque follows the optimal-$\lambda$ law
\begin{equation}
  T_g = K_g\,\Omega^2,\qquad P = T_g\Omega = K_g\,\Omega^3 ,
  \label{eq:torque_law}
\end{equation}
with $K_g$ a constant of the controller. Rotor speed and power are therefore
deterministically linked \emph{independently of the yaw angle}: a plot of
$\Omega$ against $P$ is the same curve at every $\gamma$, and ``rotor speed at a
given power'' carries no yaw information. Yaw information appears only when
$\Omega$ is compared with an external speed reference: the equilibrium of
\eqref{eq:torque_law} with the aerodynamic torque gives
$\Omega = \lambda_{\mathrm{opt}} U \cos\gamma / R$ to leading order, so
\begin{equation}
  \frac{\Omega}{U_n} = \frac{\lambda_{\mathrm{opt}}}{R}\,\frac{\cos\gamma}{h(\gamma,\lambda)}
\end{equation}
peaks at $\gamma=0$ --- but this is the C1 estimator with $P^{1/3}$ in place of
$P$ and the same anemometer, and by \eqref{eq:torque_law} it is
algebraically the same fit. The rotor-speed method is useful on fleets where the
power signal is curtailed or corrupted and $\Omega$ is not; here the reverse
holds (power is the cleanest channel), so C3 was dropped.

%======================================================================
\section{C4: yaw manoeuvres as natural experiments}\label{sec:c4}

\subsection{Idea}
Every yaw manoeuvre rotates the nacelle by a \emph{known} angle $\Delta$
(difference of \texttt{NacDir} before and after the run of updates), and a
rotation is a rotation regardless of $\delta_n$, $\delta_v$ or $b$. The power
before and after the manoeuvre is measured at two different yaw errors that
differ by a known amount, so the pair of powers samples the \emph{derivative} of
the $\cos^k$ curve at a known offset from the operating point. The vane is not
used at all; the estimate is in physical degrees. This is the only estimator in
the pipeline whose level is, in principle, free of the vane transfer function.

\subsection{Derivation}
Let a manoeuvre rotate the nacelle from $\Nt^-$ to $\Nt^+ = \Nt^- + \Delta$
(compass, clockwise positive). The controller rotates towards the wind, so
$\Delta>0$ means the wind had drifted clockwise of the nacelle. After the
manoeuvre the controller has restored its equilibrium and the yaw error is the
static misalignment, $\gamma^+ = \theta$. Before it, the wind was at the same
direction relative to the old heading plus whatever fraction of the drift had
already accumulated:
\begin{equation}
  \gamma^- = \wrap(\Wt - \Nt^-) = \theta + f\Delta,\qquad 0<f\le1 ,
  \label{eq:gamma_before}
\end{equation}
where $f=1$ if the wind shifted as a step and the nacelle followed, and $f<1$
if the wind drifted gradually and the $-6\ldots-2$~min window before the
manoeuvre already contains part of the drift. Let $P^\pm$ be the mean power in
the windows after and before, and $P_j^\pm$ the same for a reference neighbour
that did not manoeuvre. Under \eqref{eq:cosk}, with the free-stream speed
cancelling between target and reference (both in partial load, both unwaked),
\begin{equation}
  \ln R \;\equiv\; \ln\frac{P^+}{P^-} - \ln\frac{P_j^+}{P_j^-}
  \;=\; k\big[\ln\cos\theta - \ln\cos(\theta+f\Delta)\big] + \eta ,
  \label{eq:lnR}
\end{equation}
with $\eta$ the noise. Expanding $\ln\cos(\theta+x) = \ln\cos\theta - x\tan\theta
- \tfrac{x^2}{2\cos^2\theta} + O(x^3)$ with $x=f\Delta$ in radians,
\begin{equation}
  \ln R = k f \tan\theta\;\Delta \;+\; \frac{k f^2}{2\cos^2\theta}\,\Delta^2 + O(\Delta^3) + \eta .
  \label{eq:lnR_expand}
\end{equation}
\begin{proposition}[Sign and level from the manoeuvre slope]\label{prop:c4}
The coefficient of the term linear in $\Delta$ is $\beta_1 = k f\tan\theta$
(per radian). Its sign is the sign of $\theta$ in the convention
$\gamma=\Wt-\Nt$, clockwise positive, and its magnitude gives
\begin{equation}
  \hat\theta = \arctan\frac{\hat\beta_1}{k f} .
  \label{eq:c4_theta}
\end{equation}
The quadratic term is even in $\Delta$ and always positive: a manoeuvre in
either direction increases power on average, by more for larger $|\Delta|$.
Only the \emph{asymmetry} between $+\Delta$ and $-\Delta$ carries $\theta$.
\end{proposition}
Intuitively: if $\theta<0$ (rotor points clockwise of the wind), a manoeuvre with
$\Delta<0$ started from $\gamma^-=\theta+f\Delta$, even further from zero, and the
gain after it is large; a manoeuvre with $\Delta>0$ started from
$\gamma^- = \theta + f\Delta$, closer to zero or past it, and the gain is small
or negative. For $\theta=-5\dg$, $f=1$, $k=2$: $\Delta=-20\dg$ gives
$P^+/P^- = \cos^2(-5\dg)/\cos^2(-25\dg) = 1.21$, $\Delta=+20\dg$ gives
$\cos^2(-5\dg)/\cos^2(15\dg)=1.06$. This asymmetric pattern (large gain after
$-20\dg$, little after $+20\dg$) is what \texttt{PPP\_WTG17} shows in 2023.

\begin{remark}[Sign in the implementation]
\texttt{manoeuvre\_fit} returns $\theta = \arctan(-\hat\beta_1/(kf))$, i.e.\ with
the opposite sign to \eqref{eq:c4_theta}. Under the conventions of
\S\ref{sec:sensor} (both \texttt{NacDir} and \texttt{WindDir} clockwise
bearings, $\Delta$ = after minus before) the derivation above gives the positive
sign. The stored C4 levels in \texttt{cache/state\_estimates.parquet} should be
read with their sign flipped; this does not rescue C4 (see \S\ref{sec:c4_fail})
but it changes the statement in \texttt{docs/03} that C4 ``disagrees in sign''
with the consensus residual on \texttt{PPP\_WTG17}.
\end{remark}

\subsection{Estimation}
Manoeuvres are runs of \texttt{NacDir} updates less than 60\,s apart with
$|\Delta|\in[2\dg,30\dg]$ and duration below 5\,min; there are
$10^5$ per turbine, of which $9$--$26\times10^3$ have both windows entirely in
good partial-load, unwaked minutes for target and reference. The slope
$\hat\beta_1$ is obtained by a Huber regression of $\ln R$ on $\Delta$ (in
degrees, converted to radians in \eqref{eq:c4_theta}); the Huber loss
\begin{equation}
  \ell_\kappa(e) = \begin{cases} \tfrac12 e^2 & |e|\le\kappa \\ \kappa|e| - \tfrac12\kappa^2 & |e|>\kappa\end{cases}
  \label{eq:huber}
\end{equation}
is used because $\ln R$ has heavy tails (gusts inside a four-minute window).
The quadratic term is not fitted: with $\Delta$ symmetric around zero it is
orthogonal to the linear term, and it is absorbed in the intercept and residual.
The standard error is the bootstrap standard deviation over 200 resamples of the
events, and the standard error of $\hat\theta$ follows by the delta method
\begin{equation}
  \widehat{\mathrm{se}}(\hat\theta) \approx \tfrac12\big|\theta(\hat\beta_1+\mathrm{se}) - \theta(\hat\beta_1-\mathrm{se})\big| .
\end{equation}

\subsection{Why C4 is not the clean estimator hoped for}\label{sec:c4_fail}
Two systematic effects, both measured, break Proposition~\ref{prop:c4} on this
data.
\begin{enumerate}
\item \textbf{A meteorological confounder.} A \emph{null test} applies the
      manoeuvre times of \texttt{PPP\_WTG17} to the power ratio of two
      unrelated turbines that did not manoeuvre. The slope should be zero; it
      is $-1.0\pm0.2\times10^{-3}$ per degree, a third of the ``real'' slope.
      Mechanism: the sign of $\Delta$ is the sign of the wind-direction change
      (veering vs backing), and veering/backing episodes come with their own
      power evolution (frontal passages, diurnal cycle) that is not identical at
      two sites a few hundred metres apart. The reference normalisation in
      \eqref{eq:lnR} removes the common part of the speed change, not the
      direction-correlated residual. This term adds to $\beta_1$ directly and
      cannot be separated from $\theta$ without modelling it per sector.
\item \textbf{The unknown fraction $f$.} The level \eqref{eq:c4_theta} scales
      as $1/f$, and $f$ is neither known nor constant: it is the fraction of the
      wind shift that had already happened by the time of the ``before'' window.
      With gradual drift and a $\pm8\dg$ deadband, $f$ is plausibly
      $0.3$--$0.6$, i.e.\ the raw slope underestimates $|\theta|$ by a factor
      two to three.
\end{enumerate}
The reference turbines disagree among themselves by $2$--$4\dg$, and the three
training turbines are not ranked correctly. C4 is therefore kept as a feature
(the asymmetry is real and large on some turbine-years) but not used for
levels. The one open route to a label-free level is to difference the slope
against the null distribution per wind-direction sector and manoeuvre direction.

%======================================================================
\section{C5: the vane gain from manoeuvre steps, and why it is impossible here}\label{sec:c5}

If the vane responded instantaneously with gain $b$, a nacelle rotation of
$\Delta$ during steady wind would produce a vane jump of $-b\Delta$
(\eqref{eq:vm} with $\gamma\to\gamma-\Delta$), and the thousands of manoeuvres
would give $b$ directly, converting all vane-unit estimates into physical
degrees. On this data the derived vane jumps by \emph{exactly} $-\Delta$ at
every \texttt{NacDir} update. The reason is the encoding, not physics: the
stream is change-based and asynchronous, so at a \texttt{NacDir} update the
\texttt{WindDir} field still holds the value computed by the controller at the
last \texttt{WindDir} update, which was $\Nr^{\mathrm{old}} + v_m^{\mathrm{old}}$.
The forward-filled difference \texttt{WindDir}$-$\texttt{NacDir} then reads
$v_m^{\mathrm{old}} - \Delta$ until \texttt{WindDir} refreshes, whatever the
vane physically does. The apparent ``gain'' is one by construction, and no
quantity derived from the two direction channels at the 12-second scale can
recover $b$. C5 was dropped, and the implementation sets $b=1$ throughout.

%======================================================================
\section{The neighbour-consensus heading residual}\label{sec:consensus}

\subsection{Sensor algebra}
Let $c(t)$ be a consensus estimate of the true wind direction built from
neighbours (\S\ref{sec:cons_build}), $c = \Wt + \kappa$ with $\kappa$ a constant
that depends on the neighbours' own offsets. From \eqref{eq:Nr} and the
definition $\gamma=\Wt-\Nt$,
\begin{equation}
  r \;\equiv\; \wrap(\Nr - c) \;=\; \Nt + \delta_n - \Wt - \kappa \;=\; -\gamma + \delta_n - \kappa .
  \label{eq:residual}
\end{equation}
Averaging over a day, $\E[\gamma]=\theta$, so
\begin{equation}
  \boxed{\ r(t) = -\theta(t) + \delta_n(t) - \kappa\ }
  \qquad\Longrightarrow\qquad
  \Delta r = -\Delta\theta \quad\text{whenever } \Delta\delta_n = 0 .
  \label{eq:dr}
\end{equation}
\begin{proposition}
The residual identifies \emph{changes} of the static misalignment in physical
degrees with slope exactly $-1$, independent of the vane ($b$, $\delta_v$) and of
the aerodynamics, but its \emph{level} is unidentifiable: $\delta_n$ and
$\kappa$ are unknown constants.
\end{proposition}
\noindent This is the ``consensus anchor'' of the leading team. The residual is
blind to changes of $\theta$ that occur through $b$ (the wind side of
$\gamma$ does not move) --- but such changes are also invisible to the label
mechanism to first order, so they are rare in the label series.

Using $\Wr$ instead of $\Nr$ gives the same information:
$\Wr - c = \delta_n + \delta_v + (b-1)\gamma - \kappa$, which for $b=1$ is a
constant, and for $b<1$ tracks $-(1-b)\theta$, a damped version of
\eqref{eq:dr}. The implementation uses $\Nr$.

\subsection{Encoder re-references}
$\delta_n$ is piecewise constant with jumps at re-references (service, power
cycle of the encoder). A jump of $\delta_n$ moves $r$ of the affected turbine
against \emph{all} its neighbours by the same amount on the same day. The
detection rule uses exactly this signature. For every turbine $i$ and its $k=4$
nearest same-site references $j$, the daily circular median of the 10-minute
heading difference $\wrap(\Nr^i - \Nr^j)$ is unwrapped and segmented with PELT
under an $L_1$ cost (\S\ref{sec:pelt}). A level shift of at least $25\dg$ that
appears with the same sign within $\pm3$ days in a majority ($\ge60\,\%$, at
least two) of $i$'s pairs is credited to $i$'s encoder and subtracted from its
\texttt{NacDir} and \texttt{WindDir} from that date on. Sixty-seven such steps
on ten turbines were found; none on \texttt{PPP\_WTG17}. Shifts below
$25\dg$ cannot be attributed by majority (a $5\dg$ misalignment change and a
$5\dg$ re-reference look identical in every pair), which is why the state
segmentation later types boundaries with a jump above $15\dg$ as \texttt{encoder}
and refuses to carry $\Delta r$ across them.

\subsection{Building the consensus}\label{sec:cons_build}
The per-pair difference $\wrap(\Nr^i - \Nr^j)$ is not constant even with
$\theta$ and $\delta_n$ fixed: terrain steering and partial wakes make the
direction difference between two sites a function of the wind direction. The
implementation removes this by a \emph{sector correction}: within each pair,
subtract the circular median of the difference in the neighbour's $30\dg$
wind-direction sector, then take the daily median of the corrected difference
over the 10-minute bins in which both turbines operate more than 90\,\% of the
time (at least 12 bins). Days whose raw pair offset sits more than $30\dg$ from
its 31-day rolling median are masked (multi-day frame excursions shorter than
the PELT minimum size). Pairs whose daily inter-quartile range is more than
twice the turbine's median pair IQR are dropped as noisy references, and
\begin{equation}
  r_i(t) = \med_{j\in\mathcal N_i}\ \med_{\text{bins of day } t}\ \wrap\!\big(\Nr^i - \Nr^j - m_{ij}(\text{sector}_j)\big)
  \label{eq:r_def}
\end{equation}
is the daily residual, with a measured day-to-day noise of $\sigma\approx2\dg$
(MAD). The double median makes \eqref{eq:r_def} robust to one bad neighbour and
to a few bad bins; the price is that changes below $\approx1\dg$ (the two small
label steps on WTG12 and WTG14) are under the noise floor.

\subsection{Change-point segmentation (PELT)}\label{sec:pelt}
States are found by minimising, over all partitions of the day index
$1..T$ into contiguous segments $S_1,\dots,S_m$,
\begin{equation}
  \sum_{\ell=1}^{m}\ \sum_{t\in S_\ell} \big(y_t - \bar y_{S_\ell}\big)^2 \;+\; \beta\,(m-1),
  \qquad |S_\ell| \ge n_{\min},
  \label{eq:pelt}
\end{equation}
where $y_t$ is the 7-day rolling median of $r_i$ (gaps interpolated so that
states are contiguous in calendar time), $\beta=400$ and $n_{\min}=14$ days.
The pruned exact linear time algorithm solves \eqref{eq:pelt} exactly by dynamic
programming with pruning of candidate change points that can never be optimal.
The $L_2$ cost is the negative log-likelihood of a piecewise-constant Gaussian
mean, so $\beta$ plays the role of a BIC-type penalty: a break is accepted when
it reduces the residual sum of squares by more than $\beta$, which for
$\sigma_{\text{smooth}}\approx1\dg$ corresponds to a level jump $\delta$ over
$n$ days with $n\delta^2/2 \gtrsim 400$, e.g.\ $\delta\ge2.4\dg$ over 140 days
or $\delta\ge5\dg$ over 30 days. $\beta$ was chosen on the training turbines
(lower values split \texttt{PPP\_WTG17}'s 2023 into up to 13 states and
\texttt{SSS\_WTG06} into 50 --- noise, not physics). For the encoder step
detection the $L_1$ cost $\sum|y_t - \med(y_{S_\ell})|$ is used instead, so that
single outlying days do not create segments.

The state level is $\bar r_\ell = \med_{t\in S_\ell} y_t$; the first day of each
state is typed \texttt{encoder} when $|\bar r_\ell - \bar r_{\ell-1}|>15\dg$, else
\texttt{candidate}; and days within $\pm7$ of a boundary are flagged
\texttt{transition}, mirroring the scoring rule.

Measured on the training turbines against the states derived from the labels
themselves: ARI $0.47$ / $0.37$ / $0.79$ (WTG12/13/14), limited by sub-degree
label steps that $r$ cannot see; on \texttt{PPP\_WTG17} two states with the
boundary at 2024-01-06 and $\Delta r = -8.05\dg$, i.e.\ $\Delta\theta=+8\dg$.

%======================================================================
\section{The submission model}\label{sec:submission}

\subsection{Structure}
For every turbine-day the prediction is
\begin{equation}
  \hat\theta(t) = L + a + b\,x(t),\qquad
  x(t) = \bar r_{\ell(t)} - \alpha_{\sigma(\ell(t))},
  \label{eq:submission}
\end{equation}
where $\ell(t)$ is the day's state, $\sigma(\ell)$ its encoder-free segment (a
maximal run of states without an \texttt{encoder}-typed boundary), $\alpha_\sigma$
the segment anchor, $L$ the level prior, and $(a,b)$ either fixed by physics
(T0: $a=0$, $b=-1$ from \eqref{eq:dr}) or fitted on the labelled turbines (T3).
The cluster field is $\ell(t)$.

\subsection{The level prior $L$}
No label-free estimator of the level survived \S\ref{sec:c1}--\S\ref{sec:c4}.
The prior used is the \emph{vane set-point}: from \eqref{eq:theta_def} with
$b=1$ and the assumption $\delta_v=0$,
\begin{equation}
  L = s = \med_{\text{healthy days}}\ \med_{\text{operating minutes}} v_m ,
  \label{eq:L}
\end{equation}
the value at which the controller holds the vane. On the PPP site this is
$-2.7$ to $-3.7\dg$ on every turbine, on SSS $-0.8$ to $-3.3\dg$. The prior is
right when the vane is mounted true and the operator has deliberately chosen a
non-zero set-point; it is wrong by $-\delta_v$ otherwise. On the training
turbines it gives a bias of $-0.45$, $+2.5$ and $-0.8\dg$: correct within a
degree on two of three, and $4.5\dg$ off on WTG13 whose vane evidently carries
its own offset.

\subsection{The anchor $\alpha_\sigma$}
The prior \eqref{eq:L} is a statement about the turbine's \emph{typical} state.
Two choices of ``typical'' were evaluated on the training turbines:
\begin{align}
  \alpha_\sigma^{\mathrm{med}} &= \med_{t\in\sigma} \bar r_{\ell(t)}
  \quad\text{(the state holding most days),}\\
  \alpha_\sigma^{\mathrm{mean}} &= \frac{1}{|\sigma|}\sum_{t\in\sigma} \bar r_{\ell(t)}
  \quad\text{(day-weighted average of the state levels).}
\end{align}
The median anchor scores better on train (RMSE $2.18$ vs $2.48$) because the
training turbines have one dominant state. For \texttt{PPP\_WTG17}, whose two
states hold 370 and 361 days, the median anchor is a coin flip that moves the
whole prediction by $4\dg$; the mean anchor is the choice that minimises the
expected squared error under ignorance of which state is typical, and it was
submitted first. Re-anchoring per segment is what stops encoder re-references
below the $25\dg$ attribution threshold from propagating into the levels.

\subsection{T3 calibration}
On the scored training days the residual $y = \theta_{\text{label}} - L$ is
regressed on $x$ with the Huber loss \eqref{eq:huber},
\[
  (\hat a,\hat b) = \arg\min_{a,b} \sum_t \ell_{1.35}\big(y_t - a - b x_t\big),
\]
giving $\hat a=-0.73\dg$, $\hat b=-0.93$ on all three turbines: the labels
confirm the physical slope of $-1$ and shift the level by less than a degree.
Leave-one-turbine-out, the fit is \emph{worse} than T0 (RMSE $4.4$ vs $2.5$),
because only WTG13 has label steps large enough to identify $b$; holding it out
collapses $\hat b$ to zero. The T3 file is therefore the T0 file with a
$-0.7\dg$ shift and a $7\,\%$ shrink of the deltas, and the honest reading is
that three labelled turbines cannot calibrate more than the sign.

\subsection{Sign check}\label{sec:signcheck}
Two independent tests are asserted before any file is written.
\begin{enumerate}
\item The ordinary least-squares slope of the label on $x$ over the scored
      training days must be negative; measured $-0.74$. This is \eqref{eq:dr}
      seen from the labels.
\item Around 2023-01-14..18 the vane median of \texttt{PPP\_WTG13} moved by
      $-2.6\dg$ while the label moved by $+3.0\dg$. By \eqref{eq:theta_def} a
      change of the set-point $s$ (or of $\delta_v$ with the controller
      re-nulling) moves $\theta$ by the same amount in the \emph{same}
      direction, unless the label sign is opposite to the vane sign convention.
      The observed opposite motion therefore fixes: label $= -\gamma$ in the
      raw vane convention, or equivalently $\gamma$ in the convention
      ``wind minus nacelle'' after accounting for the direction in which this
      controller reports the vane. The test is asserted as ``vane and label
      move in opposite directions''; together with test~1 it pins the sign of
      the submission model, which is all that matters for scoring.
\end{enumerate}

%======================================================================
\section{Wake sectors}\label{sec:wake}
For target $i$ and each published neighbour $j$ within four times the nearest
published distance (only ratios are known; the layout is scaled), the bearing
$\beta_{ij}$ from $i$ to $j$ is exact. A minute is flagged waked when the
de-stepped wind direction lies within $\pm20\dg$ of any $\beta_{ij}$, and
all power-based estimators use unwaked minutes only. Since only 16 of the
site's turbines are published, the mask is a lower bound on real waking; residual
partial wakes are one of the sector-dependent nuisances of
Lemma~\ref{lem:tilt}.

%======================================================================
\section{Scoring rule and what it rewards}\label{sec:scoring}

\subsection{Metrics}
With $D$ the set of scored days (labelled, outside $\pm7$ days of a label step),
\begin{equation}
  \mathrm{RMSE} = \sqrt{\frac1{|D|}\sum_{t\in D}(\hat\theta_t-\theta_t)^2},\quad
  \mathrm{MAE} = \frac1{|D|}\sum_{t\in D}|\hat\theta_t-\theta_t|,\quad
  \mathrm{bias} = \frac1{|D|}\sum_{t\in D}(\hat\theta_t-\theta_t).
\end{equation}
The tie-breaker is the adjusted Rand index between the submitted clusters and
the label-derived states. For a contingency table $n_{uv}$ (cluster $u$, true
state $v$) with margins $a_u, b_v$ and $n=|D|$,
\begin{equation}
  \mathrm{ARI} = \frac{\sum_{uv}\binom{n_{uv}}{2} - \big[\sum_u\binom{a_u}{2}\sum_v\binom{b_v}{2}\big]/\binom n2}
                      {\tfrac12\big[\sum_u\binom{a_u}{2}+\sum_v\binom{b_v}{2}\big] - \big[\sum_u\binom{a_u}{2}\sum_v\binom{b_v}{2}\big]/\binom n2},
\end{equation}
which is $1$ for identical partitions up to relabelling and $0$ in expectation
for a random partition. It is invariant to cluster names and to the number of
clusters only insofar as splitting a true state costs ARI in proportion to the
split.

\subsection{Effective sample size and the paired bootstrap}
Scored days are runs of a piecewise-constant series; the error of a
piecewise-constant prediction is constant within each run, so the effective
number of independent observations is the number of states (two on
\texttt{PPP\_WTG17}, of order fifteen on the training set), not the number of
days. The organisers resample whole states with replacement, score every
submission on the same draw, and declare a submission provably better than
another only if its RMSE is lower on essentially all draws. The local harness
\texttt{yaw.evaluate.score} reproduces this with 1000 draws over
(turbine, state) blocks; the resulting 95\,\% interval for the T0 model on the
training turbines is $1.8$--$3.3\dg$ around $2.48$, which is the honest
uncertainty of every number quoted in this note.

\subsection{Implication for the strategy}
Because RMSE differences smaller than the band width are ties and ARI then
decides, and because the level of $\theta$ is what none of the physics
estimators can pin better than a few degrees, the winnable quantity is the
segmentation (ARI, from \S\ref{sec:consensus}) plus a level that is not badly
wrong. The label-free T0 model of \S\ref{sec:submission} scores $2.5\dg$ on the
training turbines against a best-constant baseline of $3.1$--$6.1\dg$; the
current T0 leader on the public board scores $3.9\dg$ with a year-static
segmentation.

%======================================================================
\section{Summary of assumptions per estimator}
\begin{center}\small
\begin{tabular}{@{}lp{3.6cm}p{2.6cm}p{2.4cm}p{4.2cm}@{}}
\toprule
 & measures & units & needs & dominant systematic \\
\midrule
C1 & peak of $P$ vs $v_m$ at fixed $U_n$ & vane ($\delta_v$, $b$) & $k$ fixed & tilt (Lemma~\ref{lem:tilt}), anemometer \\
C2 & peak of $P_i/P_j$ vs $v_m$ & vane & healthy neighbour & tilt via neighbour sector \\
C3 & $\Omega$ vs $P$ & --- & --- & degenerate under torque law \\
C4 & asymmetry of $\ln R$ in $\Delta$ & physical & $k$, $f$ & veering confounder, unknown $f$ \\
C5 & vane jump at $\Delta$ & --- & synchronous channels & impossible (encoding) \\
$r$ & $\Delta\theta$ from $\Delta r$ & physical, slope $-1$ & stable $\delta_n$ & re-references $<25\dg$ \\
$L$ & vane set-point $s$ & physical if $\delta_v=0$ & --- & $\delta_v$ (WTG13: $4.5\dg$) \\
\bottomrule
\end{tabular}
\end{center}

\end{document}
